\section{Conclusion}\label{sec:conclusion}

Quantum theory's paradoxes are often framed as demands for exotic causal structure or for an ontic wavefunction.
Following Spekkens, we instead treated many of these demands as symptoms of a category mistake: confusing causal substrate description with layer-relative inferential completion.
Six Birds Theory (SBT) supplies a language in which this separation is explicit.

In this language, record-level objecthood is determined by packaging (closure) operations induced by a record interface.
Dephasing in a pointer basis is a canonical example: dephasing is idempotent, its fixed points are record-classical states, and different record bases induce incompatible closures whose noncommutation is measurable as route mismatch.
The double slit and quantum eraser become demonstrations of when ``which-path'' is packaged into an object; the measurement problem becomes bookkeeping about enforcing a record algebra; and the cat paradox is localized as an error about when macroscopic alternatives become object-level distinctions.

We supported these claims with mechanized structural results in Lean and with fully reproducible computational experiments.
The broader message is methodological: keep causation and inference distinct, and respect a Leibnizian layer principle that forbids surplus distinctions that the layer cannot stably witness.
